

%%%%%%%%%%%%%%%
\begin{frame}[fragile]{Requêtes imbriquées}


Exemple\,: Les logements où je peux faire du ski.

\vfill

\begin{minted}[baselinestretch=.9,
               fontsize=\small,
               framesep=2mm]{sql}
select distinct nom
from  Logement, Activité
where code = codeLogement
and codeActivité = 'Ski'
\end{minted}

\vfill

C'est la syntaxe classique, ``à plat''. 

\vfill

Cas particulier\,: je ne m'intéresse qu'aux attributs de la table
\texttt{Logement}

\end{frame}

%%%%%%%%%%%%%%%
\begin{frame}[fragile]{Avec quantification existentielle}

Formulation légèrement différente\,: les logements où \alert{il existe} 
une activité 'Ski'

\vfill

\begin{minted}[baselinestretch=.9,
               fontsize=\small,
               framesep=2mm]{sql}
select nom
from  Logement
where exists (select '' from Activité
       where code = codeLogement
       and codeActivité = 'Ski')
\end{minted}

\vfill

Note\,: la clause \sqlkw{select} de
la requête imbriquée ne sert à rien. Je m'intéresse 
 uniquement à l'existence
(ou non) des nuplets de \texttt{Activité}.
\end{frame}

%%%%%%%%%%%%%%%
\begin{frame}[fragile]{Avec condition d'appartenance}

Autre formulation\,:
Les logements qui font partie de ceux où on trouve une activité 'Ski'

\vfill

\begin{minted}[baselinestretch=.9,
               fontsize=\small,
               framesep=2mm]{sql}
select nom
from  Logement
where code in  (select codeLogement from Activité
                where codeActivité = 'Ski')
\end{minted}

\vfill

Note\,: correspondance entre les attributs \texttt{code}
du bloc principal et \texttt{codeLogement} de
la requête imbriquée

\vfill
\alert{Requêtes dites corrélées} à cause
du lien établi entre les deux blocs.
\end{frame}


%%%%%%%%%%%%%%%
\begin{frame}[fragile]{Quand utiliser l'imbrication\,?}

\alert{Condition}\,: c'est une jointure, mais on construit le résultat avec les nuplets d'une seule
relation. Dans l'algèbre, c'est une \alert{semi-jointure}

\vfill

L'imbrication n'apporte pas d'expressivité supplémentaire. Le choix est une question de goût.

\vfill

Critères à prendre en compte\,:

\vfill

\begin{redbullet}
  \item La lisibilité compte\,! Au-delà d'un niveau d'imbrication,
      l'interprétation devient laborieuse
  \item On ne peut pas construire le nuplet-résultat avec les nuplets des sous-requêtes
  
\end{redbullet}

\end{frame}


%%%%%%%%%%%%%%%
\begin{frame}[fragile]{Critère de lisibilité\,: un exemple (à étudier)}

Les voyageurs qui sont allés en Corse.

\vfill

\begin{minted}[baselinestretch=.9,
               fontsize=\small,
               framesep=2mm]{sql}
select prénom, nom
from  Voyageur as V
where  exists (select '' from Séjour as S
               where V.idVoyageur = S.idVoyageur
               and exists (select '' from Logement
                  where codeLogement=code
                  and lieu='Corse') )
\end{minted}
\vfill
La même requête, à plat.

\begin{minted}[baselinestretch=.9,
               fontsize=\small,
               framesep=2mm]{sql}
select distinct V.prénom, V.nom
from  Voyageur as V, Séjour as S, Logement as L
where V.idVoyageur = S.idVoyageur
and   codeLogement=code
and lieu='Corse'
\end{minted}

\end{frame}

\begin{frame}{À retenir}

Tant qu'on n'exprime pas de \alert{négation}, l'imbrication
n'apporte que des manières alternatives d'exprimer une même requête.
\vfill

\begin{redbullet}
 \item Bien comprendre les équivalences (étudier les exemples)
  \item La lisibilité est le principal critère\,: SQL, c'est apprendre
    à  exprimer des requêtes avec élégance\,!
  \item Une requête imbriquée évite naturellement des doublons
      qu'il faut éliminer dans la version ``à plat''
\end{redbullet}

\vfill

Une séquence pour rien\,? Non car l'imbrication est indispensable
pour les \alert{négations}. À suivre.
\end{frame}

